% !TEX program = xelatex
\documentclass[a4paper]{article}
\usepackage{amsthm}
\usepackage{amssymb}
\usepackage{bm}
\usepackage{mathtools}
\usepackage[x11names]{xcolor}
\usepackage{xparse}
\usepackage{fontspec}
\usepackage{unicode-math}
\setromanfont{DovesType-Regular.otf}
\setsansfont{Andika}
\setmathfont{Asana Math}[Scale=1]

% \usepackage{pstricks}
\usepackage{varwidth}
\usepackage{siunitx}
\usepackage{graphicx}
\usepackage[margin=1cm]{geometry}
\usepackage[most]{tcolorbox}
\usepackage{pgfplots}
\pgfplotsset{compat=newest}
\tcbuselibrary{skins,xparse,poster,breakable}
% \usetikzlibrary{fadings}
\usetikzlibrary{calc, plotmarks, shapes, shapes.geometric, positioning, angles, intersections, quotes, through, patterns, turtle, arrows.meta}
\usetikzlibrary{decorations.markings}
% \usepackage{etoolbox}
% \usepackage{tkz-euclide}
% \usepackage{xlop}
% \newcommand\hole[2]{#1}  % for use with xlop

%%%%%%%%%%%%%%%%%%%%%%%%%%%%%%%%%%%%%%%%%%%%%%%%%%%%%%%%%
\newcommand\markangle[9]{% origin X Y radius radiusmark mark colour opacity
%  % fill red circle offset-from-centre
  \begin{scope}
    \path[clip] (#1) -- (#2) -- (#3);
    \fill[color=#7,fill opacity=#8,draw=black,name path global=pcircle]  % global declaration required otherwise pcircle is not seen by the `named intersections=' lines below.
    (#1) circle (#4);
  \end{scope}
  % middle calculation
  \path[name path=line one] (#1) -- (#2);
  \path[name path=line two] (#1) -- (#3);
  \path[%
  name intersections={of=line one and pcircle, by={inter one}},
  name intersections={of=line two and pcircle, by={inter two}}
  ] (inter one) -- (inter two) coordinate[pos=#9] (place);
  % put mark
  \node at ($(#1)!#5!(place)$) {\scriptsize{#6}};
}

% \newcommand{\condSoln}[2]{\ifcsdef{r@#1}{#2}{}}

% \newcommand\fadingtext[3][]{%
%    \begin{tikzfadingfrompicture}[name=fading letter]
%      \node[text=transparent!0,inner xsep=0pt,outer xsep=0pt,#1] {#3};
%    \end{tikzfadingfrompicture}%
%    \begin{tikzpicture}[baseline=(textnode.base)]
%      \node[inner sep=0pt,outer sep=0pt,#1](textnode){\phantom{#3}};
%      \shade[path fading=fading letter,#2,fit fading=false]
%      (textnode.south west) rectangle (textnode.north east);%
%    \end{tikzpicture}%
% }

\definecolor{JISpurple}{RGB}{89,72,122}
\definecolor{JISivory}{RGB}{241,234,221}
\definecolor{JIStaupe}{RGB}{183,156,154}
\definecolor{PaleGreen}{RGB}{240,255,240} % 'Honeydew'

\AddToHook{shipout/background}{%
    \put (0in,-\paperheight){\includegraphics[width=\paperwidth,height=\paperheight]{images/R10V5-R.png}}%
}

\newcommand\numberthis{\addtocounter{equation}{1}\tag{\theequation}}

\newtcolorbox{MyOuterBox}{%
  enhanced,
  breakable,
  frame style=JISpurple,
  colback=JISivory,
  colframe=JISpurple,
  title={\includegraphics[width=0.9cm,height=0.9cm]{images/JIS Final Logo FA-02.png}\raisebox{3mm}{\Large{Maths Challenge}\hspace{26em} \Large{\bfseries\sffamily 10}}},
}

\newtcolorbox{MyInnerBox}[2][]{enhanced,%empty,
coltitle=JISpurple,colback=white,
fonttitle=\bfseries\sffamily,
attach boxed title to top left={yshift=-1.5mm},
boxed title style={empty, size=small, top=1mm, bottom=0pt},
varwidth boxed title=0.5\linewidth,
frame code={
  \path (title.east|-frame.north) coordinate (aux);
\path[draw=JISpurple, line width=0.5mm, rounded corners,fill=white]
(frame.west) |- ([xshift=-2.5mm]title.north east) to[out=0, in=180] ([xshift=7.5mm]aux)-|(frame.east)|-(frame.south)-|cycle;
},
title={#2},#1}

\newtcolorbox{MyInnerSplitBox}[2][]{enhanced,%empty,
bicolor,sidebyside,sidebyside align=top seam,
righthand width=8.5cm,colbacklower=white,
sidebyside gap=5mm,
coltitle=JISpurple,colback=white,
fonttitle=\bfseries\sffamily,
attach boxed title to top left={yshift=-1.5mm},
boxed title style={empty, size=small, top=1mm, bottom=0pt},
varwidth boxed title=0.5\linewidth,
frame code={
  \path (title.east|-frame.north) coordinate (aux);
\path[draw=JISpurple, line width=0.5mm, rounded corners,fill=white]
(frame.west) |- ([xshift=-2.5mm]title.north east) to[out=0, in=180] ([xshift=7.5mm]aux)-|(frame.east)|-(frame.south)-|cycle;
},
title={#2},#1}


\newtcolorbox{MySolutionBox}{%
  enhanced,
  breakable,
  frame style=JISpurple,
  colback=PaleGreen, colframe=green,
  title={\Large Solution},
  drop fuzzy shadow,
  halign=left,
}

%%%%%%%%%%%%%%%%%%%%%%%%%%%%%%%%%%%%%%%%%%%%%%%%%%
\newtoggle{SOLUTION}
%%% Uncomment the appropriate line below to show solutions %%%
% \toggletrue{SOLUTION}
\togglefalse{SOLUTION}
%%%%%%%%%%%%%%%%%%%%%%%%%%%%%%%%%%%%%%%%%%%%%%%%%


%%%%%%%%%%%%%%%%%%%%%%%%%%%%%%%%%%%%%%%%%%%%%%%%%%
%%%%%%            DOCUMENT BEGINS           %%%%%%
%%%%%%%%%%%%%%%%%%%%%%%%%%%%%%%%%%%%%%%%%%%%%%%%%%
\begin{document}


  \begin{MyOuterBox}
    \iftoggle{SOLUTION}{Here are the full, or partial solutions.
    }{
      Welcome to this week's Maths Challenge!\\
      Solutions must be explained in detail, responses with just the answer may be ignored.\\
      Drop your solution in the box in the staffroom by Tuesday.
    }
       \begin{MyInnerBox}{Year 8 and below}
      Some molten rhenium is being poured into three containers, A, B and C. Container A gets 1 litre plus one third of what is left. Container B then receives 6 litres plus one third of what remains, then Container C gets the rest which is 40 litres. What is the volume of rhenium poured into Container B?
      \iftoggle{SOLUTION}{%conditional output begin
      \begin{MySolutionBox}
        Let \(R\) be the total volume of rhenium in litres.\par
        \begin{tabular}{c|c|c|c|c|l}
          Step & A & B & C & Remainder&\\\hline
          1 & \(1 + \frac{R-1}{3}\) & ? & \(40\) & \(R-\left(1+\frac{R-1}{3}\right)\) & Simplify the remainder,\\[3mm]
            & & & & \(\frac{2R-2}{3}\) &  simplify A and work out B in Step 2:\\
          2 & \(\frac{R+2}{3}\) & \(6+\frac{\frac{2R-2}{3}-6}{3}\) & \(40\) & \(0\) & Simplify:\\[3mm]
          3 & \(\frac{R+2}{3}\) & \(6+\frac{2R-20}{9}\) & \(40\) & \(0\) & Now we can equate the sum of A, B \& C to R.\\
        \end{tabular}
        \begin{align*}
          R &= \frac{R+2}{3} + 6 + \frac{2R-20}{9} + 40\\
          R &= \frac{3R+6}{9} + \frac{2R-20}{9} + 46\\
          R &= \frac{5R-14}{9} + 46\\
          9R - 5R &= -14 + 414\\
          4R &= 400\\
          R &= \SI{100}{\litre}
        \end{align*}
        Alternatively, we can attack the problem from the end, begin with Container C with \(40\) and work backwards:\par
        We know that C is \(40\) and from the question, this is \(\frac{2}{3}\) of the quantity that B got \(\frac{1}{3}\) of. So \(40\times\frac{1}{2}=20\).\par
        But we know that B got this plus \(6\), so B got \(26\). Then for Container A: it has \(\frac{1}{2}\) of \(66\), (Container B and Container C's volume) plus \(1\) more litre, that is \(34\) litres. There is less algebra this way but perhaps it is trickier to think about.\par
        So the total volume of rhenium is \(34 + 26 + 40 = \SI{100}{\litre}\).
      \end{MySolutionBox}
    }{}%conditional output end
    \end{MyInnerBox}


    \vspace{0.4cm}
          \begin{MyInnerSplitBox}{Year 8 and above}
        The diagram contains four squares. The smallest square has an area of \SI{5}{\square\cm}.
Find the area of the shaded triangle.
      \iftoggle{SOLUTION}{%conditional output begin
        \begin{MySolutionBox}
          If you remember that the area of a triangle is half the base times the height, this problem becomes easier. Can you see that the dashed line EC is parallel to the side of the triangle JK? This is because they are both the diagonals in squares so their slope must be the same.\par
          If we move the vertex of the blue triangle at E along the dotted line to C, the base JK stays the same and so does its height, thus its area also stays the same. In other words, we have shown that \(\bigtriangleup KEJ\) has the same area as the green \(\bigtriangleup CJK\).\par
          Now we can do the same trick again. Line BK is parallel to line JC, so if we now move vertex K of the green triangle along the dotted line BK to B, we end up with \(\bigtriangleup BCJ\).\par
          The triangle fills half of the square, which has area four times that of the smallest square. So the triangle has an area of \SI{10}{\square\cm}.
         \end{MySolutionBox}
       }{}%conditional output end
      \tcblower
      \begin{tikzpicture}[scale=1]
        \coordinate (A) at (0,0);
        \coordinate (B) at (1,0);
        \coordinate (C) at (3,0);
        \coordinate (D) at (8,0);
        \coordinate (E) at (8,5);
        \coordinate (F) at (3,5);
        \coordinate (G) at (3,2);
        \coordinate (H) at (0,2);
        \coordinate (I) at (1,1);
        \coordinate (J) at (1,2);
        \coordinate (K) at (0,1);
        \draw (A) -- (D) -- (E) -- (F) -- (G) -- (H) -- cycle;
        \draw (B) -- (J);
        \draw (K) -- (I);
        \draw (C) -- (G);
        \filldraw[fill=LightSkyBlue1,opacity=0.5] (K) -- (J) -- (E) --cycle;
        \node at (0.5,0.5) {\(5\)};
        \iftoggle{SOLUTION}{
          \draw[dashed,-Stealth] (E) -- (C);
          \draw[dotted,-Stealth] (K) -- (B);
          \node[below] at (A) {A};
          \node[below] at (B) {B};
          \node[below] at (C) {C};
          \node[above] at (E) {E};
          \node[left] at (K) {K};
          \node[above] at (J) {J};
          \filldraw[fill=DarkOliveGreen1,opacity=0.5] (K) -- (J) -- (C);
        }{}
      \end{tikzpicture}
    \end{MyInnerSplitBox}


  \end{MyOuterBox}

%%%%%%%%%%%%%%%%%%%%%%%%%%%%%%%%%%%%%%%%%%%%%%%%%%
\end{document}



